\documentclass[12pt]{article}
\usepackage[margin=1in]{geometry}
\usepackage[T1]{fontenc}
\usepackage{xcolor}
\usepackage{soul}
\definecolor{garnet}{HTML}{73000A}

% Highlight commands using soul (supports line breaks)
% Yellow = filler
\newcommand{\ylw}[1]{{\sethlcolor{yellow}\hl{#1}}}
% Red = unsupported claim
\newcommand{\rd}[1]{{\sethlcolor{red!30}\hl{#1}}}
% Blue = too much technical detail
\newcommand{\blu}[1]{{\sethlcolor{cyan!30}\hl{#1}}}

\setlength{\parindent}{0pt}
\setlength{\parskip}{0.3em}

\begin{document}

\noindent
{\large\textbf{REVIEW KEY:}}\\
\ylw{Yellow} = Filler (words without benefit) \quad
\rd{Red} = Unsupported claim \quad
\blu{Blue} = Too much technical detail, not enough about student

\vspace{0.5em}
\noindent
{\color{garnet}\rule{\textwidth}{0.5pt}}
\vspace{0.5em}

\noindent
Department of Mechanical Engineering\\
University of South Carolina\\
Columbia, SC 29208\\
\today

\vspace{0.5em}

\noindent
SC Space Grant Consortium\\
College of Charleston\\
Charleston, SC 29424

\vspace{0.5em}

\noindent
\textbf{Re: Faculty Recommendation for Mateo Altomare --- Undergraduate Research Award}

\noindent
{\color{garnet}\rule{\textwidth}{0.5pt}}

\vspace{0.3em}

\noindent
Dear Members of the Evaluation Committee:

I recommend Mateo Altomare for the SC Space Grant Undergraduate Research Award \ylw{without reservation}. As his faculty advisor in the Integrated Multiphysics and Systems Engineering Laboratory (iMSEL) since June 2025, I have directly supervised his research on multi-spectral sensor fusion for autonomous maritime perception and can speak specifically to his capabilities and readiness for the proposed project.

Mateo joined iMSEL as a rising sophomore with no prior experience in machine learning or computer vision. \rd{Within his first month}, he progressed from introductory Python to implementing camera calibration routines for our thermal-visible sensor rig. \rd{By his third month}, he was independently writing PyTorch training scripts for our Detection Transformer (DETR) model on maritime imagery, having \rd{taught himself} transfer learning from \blu{RGB-pretrained backbones to thermal feature extraction}. He is now co-first-author on a paper submitted to ASME IDETC 2026 presenting \blu{a depth-aware linear affine transformation method for pixel-level IR-to-RGB registration}---a contribution that required him to understand \blu{homography estimation, parallax geometry, and depth-conditioned coordinate transforms} \rd{at a level I typically see in second-year graduate students}.

Three specific accomplishments illustrate his technical range. First, he deployed real-time YOLO object detection on a Raspberry Pi for automated target identification in \blu{marine radar Plan Position Indicator (PPI) displays, which required him to optimize inference for a platform with roughly 4\,GB of RAM and no GPU}. Second, he engineered our \blu{4-channel multi-spectral detection pipeline that fuses thermal IR and RGB imagery using scene-based keypoint extraction for cross-modal calibration}---this pipeline now generates all training data for his proposed LDAR framework. Third, he \rd{systematically} curated and annotated approximately 750 thermal-visible frame pairs spanning 1--75\,m depth range, \blu{understanding that the quality of registration ground truth directly constrains the accuracy of any learned model}.

His proposed research---Learned Depth-Adaptive Registration (LDAR)---addresses a real limitation in our current workflow. \blu{Our existing method handles parallax by calibrating separate homographies for foreground objects at 0.5\,m depth intervals, then applying them via manual segmentation masks. This works but requires knowing object depth in advance. Mateo's proposal to train a MobileNetV3-based monocular depth estimator from thermal imagery, then condition a homography network on that estimated depth, would eliminate the need for manual depth specification and LiDAR hardware.} I have reviewed his proposed architecture and training plan in detail; the approach is technically sound, the 750-pair dataset is sufficient for fine-tuning with transfer learning, and the sub-3-pixel accuracy target is ambitious but achievable given our preliminary results with fixed-depth homographies.

This work has direct relevance to NASA's Artemis program. Lunar South Pole rovers will navigate permanently shadowed regions where visible cameras are ineffective and GPS does not exist. \blu{Thermal imaging exploits the 25--50\,K temperature variations measured by the Diviner Lunar Radiometer in these regions, but precise thermal-visible registration is required for obstacle detection.} The depth-adaptive fusion techniques Mateo is developing address exactly this sensor alignment problem for platforms where every gram of hardware matters.

Beyond research, Mateo co-leads the Gamecock Motorsports Formula SAE chassis team and maintains a 3.91\,GPA. He is \rd{the kind of student who arrives early to lab}, \rd{keeps a detailed research notebook}, and \rd{follows up on results that don't match expectations rather than discarding them}. I am \ylw{confident} he will \rd{complete the proposed 350+ hours of research productively} and \rd{deliver the stated outcomes}.

I recommend Mateo Altomare \ylw{strongly} and \ylw{without reservation}.

\vspace{0.8em}

\noindent
Sincerely,

\vspace{1.5em}

\noindent
Dr.\ Yi Wang\\
Associate Professor\\
Department of Mechanical Engineering\\
University of South Carolina\\
\textit{ywang@sc.edu}

\newpage

\section*{REVIEW SUMMARY}

\subsection*{\ylw{YELLOW} --- Filler / Padding}

\begin{enumerate}
\item \textbf{``without reservation'' (para 1):} Generic phrase every strong letter uses. Adds no information.
\item \textbf{``confident'' (closing para):} Weak intensifier. Replace with evidence: ``Based on his completion of all Fall 2025 milestones ahead of schedule...''
\item \textbf{``strongly'' (final line):} Redundant intensifier when paired with ``without reservation.''
\item \textbf{``without reservation'' (final line):} Second use. Repetition weakens the letter.
\end{enumerate}

\subsection*{\rd{RED} --- Unsupported Claims}

\begin{enumerate}
\item \textbf{``Within his first month'' / ``By his third month'':} Timeline claims need evidence (git commits, dated lab notebooks, milestone logs). Otherwise soften or remove.
\item \textbf{``taught himself'':} How do you know? More accurate: ``mastered through independent study'' with evidence.
\item \textbf{``at a level I typically see in second-year graduate students'':} Subjective comparison without calibration. Specify what makes it graduate-level.
\item \textbf{``systematically'':} Generic adverb. Show the system he used.
\item \textbf{``arrives early to lab'':} Unverifiable without attendance data.
\item \textbf{``keeps a detailed research notebook'':} Have you reviewed it? Describe what makes it detailed.
\item \textbf{``follows up on results that don't match expectations'':} Give ONE specific example with dates.
\item \textbf{``complete 350+ hours productively'' / ``deliver stated outcomes'':} Based on what past performance data?
\end{enumerate}

\subsection*{\blu{BLUE} --- Too Much Technical Detail}

\begin{enumerate}
\item \textbf{Para 2 technical terms:} ``RGB-pretrained backbones,'' ``homography estimation, parallax geometry, depth-conditioned coordinate transforms'' describe the MATH, not Mateo. Reframe around what he learned/demonstrated.
\item \textbf{Para 3 accomplishment details:} PPI displays, 4\,GB RAM, 4-channel pipeline, keypoint extraction---bury the student's initiative under jargon. Lead with what Mateo \emph{did}, then briefly name the technology.
\item \textbf{Para 4 (nearly all):} Reads like a technical proposal review. The committee doesn't need MobileNetV3 architecture details. Compress to: ``His proposed method solves a real limitation. I've reviewed it and find it technically sound and achievable.''
\item \textbf{Para 5 Diviner details:} Reviewers know the lunar south pole is challenging. The 25--50\,K Diviner temperature data is overkill for a rec letter. One sentence of NASA context suffices.
\end{enumerate}

\subsection*{OVERALL ASSESSMENT}

\textbf{Strengths:} Specific numbers (750 frames, 3.91 GPA), verifiable outputs (IDETC paper, Formula SAE role), clear narrative arc (no experience $\to$ co-author).

\textbf{Fix:} Cut technical detail by 50\%. Replace unsupported behavioral claims with one concrete anecdote. Remove all filler phrases. Every sentence should answer: ``What does this show about \emph{Mateo}?''

\end{document}
